\documentclass[11pt]{amsart}
\ifdefined\pdfinfoomitdate\pdfinfoomitdate=1\fi
\ifdefined\pdftrailerid\pdftrailerid{}\fi
\ifdefined\pdfsuppressptexinfo\pdfsuppressptexinfo=-1\fi
\usepackage[T1]{fontenc}
\usepackage{lmodern}
\usepackage{microtype}
\usepackage{amsmath,amssymb,mathtools}
\usepackage{enumitem}
\usepackage{xcolor}
\usepackage[pdfusetitle,hidelinks]{hyperref}
\hypersetup{
  pdftitle={Generic Exact Reconstruction of Reflected-Cap Data from an Erased Flat Boundary},
  pdfauthor={Oleksiy Babanskyy},
  pdfsubject={Exact inverse reconstruction for proper-zero reflected-cap data},
  pdfkeywords={convex polytopes, inverse geometry, Voronoi reconstruction, exact computation}
}

\newtheorem{theorem}{Theorem}[section]
\newtheorem{proposition}[theorem]{Proposition}
\newtheorem{lemma}[theorem]{Lemma}
\newtheorem{corollary}[theorem]{Corollary}
\theoremstyle{definition}
\newtheorem{definition}[theorem]{Definition}
\newtheorem{remark}[theorem]{Remark}

\newcommand{\R}{\mathbb{R}}
\newcommand{\conv}{\operatorname{conv}}
\newcommand{\aff}{\operatorname{aff}}
\newcommand{\relint}{\operatorname{relint}}
\newcommand{\vol}{\operatorname{vol}}
\newcommand{\dist}{\operatorname{dist}}
\newcommand{\GammaM}{\Gamma(M_J)}

\title[Reconstruction from an erased flat boundary]{Generic exact reconstruction of reflected-cap data\\from an erased flat boundary}
\author{Oleksiy Babanskyy}
\address{Independent researcher}
\email{aconsciousfractal@users.noreply.github.com}
\urladdr{https://orcid.org/0009-0001-6176-6208}
\date{July 31, 2026}

\raggedbottom

\begin{document}

\begin{abstract}
Let a convex three-polytope in a hyperplane of $\R^4$ carry one reflected
pyramidal cap on each facet, with the caps independently rotated about their
facets.  We study the proper-zero regime: at least one cap remains nonflat and
at least one cap returns to the core hyperplane.  The flat caps then merge with
the core and erase their internal source-layer boundaries.

We define an intrinsic boundary-germ skeleton of the merged flat component.
Its vertices satisfying a perpendicular-bisector neighbour test are exactly
the hidden reflected apexes, provided no core vertex satisfies the same test.
This yields exact reconstruction and uniqueness of every compatible
decomposition.  For each fixed realized core and visible-facet mask, centres
violating the hypothesis lie in a finite union of proper spheres and quadrics,
and hence form a measure-zero set.

Beyond the generic statement we prove two unconditional structure results.
Uniqueness holds whenever every core vertex that passes the test lies on a
visible hinge facet.  Moreover, any two distinct compatible decompositions
must differ in at least five reflected sites on each side; consequently a
proper-zero datum with at most four hidden facets is unique without any
genericity hypothesis.  The five-site bound rests on a nonexistence theorem
for an eight-point spherical configuration, proved by four polynomial
certificate identities that are re-verified by exact expansion, together with
an elementary positivity analysis.

We also give an exact rational reference algorithm, a certifying variant
whose two comparison conditions are each provably necessary, and a complete
16,744-instance finite replay.  The finite replay supports the
implementation; it is not used as the proof.
\end{abstract}

\maketitle

\section{Introduction}

Inverse Voronoi and Dirichlet problems usually begin with a supplied
tessellation, cell complex, embedded diagram, or collection of cell
boundaries.  The input considered here is coarser.  Several reflected sites
and their facet boundaries have been erased by a coplanar union before the
inverse problem begins.

The model is elementary.  A convex polytope $P$ lies in a three-dimensional
hyperplane $H\subset\R^4$, an interior point $o$ is fixed, and the reflection
of $o$ in each facet plane is used as the apex of a pyramid.  Each pyramid is
then rotated through an independent angle about its base facet.  When an angle
is zero, the corresponding pyramid is coplanar with $P$.  All such flat
pyramids merge with $P$ into one three-dimensional subset $M_J\subset H$.
Only the raw union is observed; cell labels and internal flat seams are not.

The central issue is therefore not the forward reflection construction.  That
construction is classical and appears explicitly in Ash and Bolker's
Theorem~2~\cite{AshBolker1985}.  The issue is whether the hidden reflected
sites can be identified from the exposed boundary after the internal seams
have disappeared.

A naive graph is insufficient.  Taking the polygonal one-skeleton after every
coplanar merge can swallow a collinear point at which the local boundary germ
changes.  Retaining the sides of the original pyramids repairs that example but
illegitimately remembers the unknown decomposition.  We instead use the
noncoplanar crease graph of the maximal exposed planar sheet germs.  It keeps a
collinear point exactly when the local incident-sheet germ changes, and removes
it when it is only a presentation subdivision.

Our main result is conditional but size-uniform.

\begin{theorem}[informal]
In dimension three, suppose a reflected-cap datum is proper-zero and no core
vertex passes the intrinsic bisector-neighbour test.  Then the passing vertices
are precisely the hidden reflected apexes, the core is recovered by their
bisector halfspaces together with the visible hinge halfspaces, and the datum
is unique up to ambient Euclidean congruence.  For a fixed realized core and
visible mask, the excluded centres form a measure-zero subset of the core.
\end{theorem}

At excluded centres the conditional theorem is silent, but the ambiguity it
leaves open is strongly constrained.

\begin{theorem}[informal]
Uniqueness already holds whenever every core vertex that passes the test lies
on at least one visible hinge facet.  In general, two distinct compatible
decompositions of one proper-zero raw union must differ in at least five
reflected sites on each side.  Hence a proper-zero datum with at most four
hidden facets is unique, with no hypothesis on its centre.
\end{theorem}

The five-site bound reduces, through an exchange graph and an elementary
incidence count, to the nonexistence of a specific eight-point spherical
configuration.
That nonexistence is proved by a computer-assisted but fully verifiable
argument: four explicit polynomial identities, found by a Gr\"obner search and
re-verified independently by exact expansion, force the configuration into a
one-parameter family that an elementary positivity analysis then empties.
Section~\ref{sec:exchange} contains the complete argument.

The theorems do not assert uniqueness for all data, a result in higher
dimension, or novelty or priority relative to all possible prior art.  Those
boundaries are repeated in Section~\ref{sec:limitations}.

\section{The reflected-cap model and exact contacts}
\label{sec:model}

Let $H\cong\R^3$ be a fixed hyperplane in $\R^4$, and let $e$ be a unit normal
to $H$.  Let $P\subset H$ be a full-dimensional convex polytope and
$o\in\operatorname{int}P$.  Translate $o$ to the origin and write the
irredundant unit-normal description
\[
 P=\{x\in H:n_i\cdot x\le h_i,\ i=1,\ldots,m\},
 \qquad \|n_i\|=1,\quad h_i>0,
\]
with facets $F_i=P\cap\{n_i\cdot x=h_i\}$.
For $\theta_i\in(-\pi,\pi)$ set
\[
 a_i(\theta_i)=h_i(1+\cos\theta_i)n_i+h_i\sin\theta_i e,
 \qquad
 C_i(\theta_i)=\conv(F_i\cup\{a_i(\theta_i)\}).
\]
The excluded value $\pi$ collapses the apex to $o$.  At $\theta_i=0$,
$a_i=2h_i n_i$, the reflection of $o$ in $\aff F_i$: at zero angle the sites
are exactly those of Ash and Bolker's forward
construction~\cite{AshBolker1985}.

Two data are called congruent when an ambient Euclidean isometry carries the
core, centre, and unordered cap family of one to those of the other.  Thus the
angle coordinates are used only to parametrise the geometric caps.

The datum is $D=(P,o,\theta)$ and its raw union is
\[
 X(D)=P\cup\bigcup_{i=1}^m C_i(\theta_i)\subset\R^4,
\]
retained as a bare Euclidean subset.  Put
\[
 J=\{i:\theta_i=0\},\qquad I=[m]\setminus J.
\]
The datum is \emph{proper-zero} when $I$ and $J$ are both nonempty.  Its
merged flat component is
\[
 M_J=P\cup\bigcup_{j\in J}C_j(0)\subset H.
\]

\begin{proposition}[exact contacts]
\label{prop:contacts}
For every angle vector $\theta\in(-\pi,\pi)^m$,
\[
 P\cap C_i(\theta_i)=F_i,
 \qquad
 C_i(\theta_i)\cap C_j(\theta_j)=F_i\cap F_j\quad(i\ne j).
\]
Thus the cells form an embedded face-to-face polyhedral complex whose contact
poset is independent of the angles.
\end{proposition}

\begin{proof}
Use the Voronoi diagram of the sites
$0,a_1(\theta_1),\ldots,a_m(\theta_m)$.  The sites are distinct: equality of
two nonzero $H$-components would force two irredundant outward facet normals
to agree.  For $x\in P$,
\begin{align*}
 \|x-a_i\|^2-\|x\|^2
 &=\|a_i\|^2-2x\cdot a_i\\
 &=2h_i(1+\cos\theta_i)(h_i-n_i\cdot x)\ge0.
\end{align*}
The first factor is positive on $(-\pi,\pi)$.  Hence $P$ lies in the Voronoi
cell $V_0$ of $0$, and $F_i\subset V_0\cap V_i$.  The apex $a_i$ is in the
interior of $V_i$.  Convexity of $V_i$ gives
$C_i\setminus F_i\subset\operatorname{int}V_i$.  Therefore no point of
$P\cap C_i$ can lie outside $F_i$.

If a point of $C_i\cap C_j$ lay outside one of the two hinge facets, the same
argument would place it in one open Voronoi cell and simultaneously in a
different open cell or in $V_0$.  This is impossible.  The reverse inclusion
$F_i\cap F_j\subset C_i\cap C_j$ is immediate.
\end{proof}

\section{Intrinsic recovery in the proper-zero regime}
\label{sec:intrinsic-recovery}

For a three-dimensional polyhedral subset $Y\subset\R^4$, let
$\operatorname{FlatReg}_3(Y)$ be the set of points having a neighbourhood in
$Y$ that is relatively open in one affine three-plane.  This set is determined
by $Y$ alone.

\begin{proposition}[proper-zero intrinsic recovery]
\label{prop:intrinsic-recovery}
If $I\ne\varnothing$, the raw subset $X(D)$ intrinsically determines $M_J$,
every nonflat cap and its hinge facet, the centre $o$, and the fully reflected
site associated with every nonflat hinge.
\end{proposition}

\begin{proof}
The closures of the connected components of
$\operatorname{FlatReg}_3(X)$ are $M_J$ and the nonflat caps.  Indeed, the
zero-angle cells join the core through facets inside $H$ and form one
three-dimensional flat component, while every nonflat cap has a different
affine hull.  Proposition~\ref{prop:contacts} accounts for every contact.
The codimension-one adjacency graph of these closures is a star.

If $|I|\ge2$, the unique star centre is $M_J$.  If $|I|=1$, let $i$ be the
nonflat index, let $V=\vol_3(P)$, and let
$v_k=\vol_3\conv(\{o\}\cup F_k)$.  The radial pyramids tile $P$, so
$\sum_kv_k=V$.  Rotation preserves the intrinsic volume of a cap, and hence
\[
 \vol_3(M_J)=2V-v_i>v_i=\vol_3(C_i).
\]
Thus in the two-component case the larger component is $M_J$.

Fix one recovered nonflat cap, with apex $A_i$ and hinge $F_i$.  Let $p_i$ be
the orthogonal projection of $A_i$ to $\aff F_i$, let
$r_i=\dist(A_i,\aff F_i)$, and let $\nu_i$ be either unit normal to
$\aff F_i$ inside $H$.  Rotation about $\aff F_i$ preserves the radius in its
two-dimensional normal plane, so
\[
 \{p_i-r_i\nu_i,p_i+r_i\nu_i\}=\{o,a_i(0)\}.
\]
The centre belongs to the interior of $M_J$.  The other point does not belong
to $M_J$: apply Proposition~\ref{prop:contacts} to the angle vector in
which cap $i$ is also set to zero.  Its apex meets neither the core nor another
cap away from prescribed lower-dimensional contacts.  Consequently $o$ is the
unique one of the two candidates in $\operatorname{int}M_J$.  Reflection of
$o$ in the recovered hinge then gives the fully reflected site.
\end{proof}

Every compatible decomposition of a given proper-zero raw union therefore has
the same $M_J$, centre, nonflat caps, and visible hinge facets.

\section{The intrinsic boundary-germ skeleton}
\label{sec:skeleton}

Let an \emph{exposed planar sheet} of $\partial M_J$ be the closure of a
maximal connected planar piece of its relative-regular locus.  Internal
coplanar subdivisions are not retained.

\begin{definition}[intrinsic germ skeleton]
\label{def:skeleton}
Atomise the boundaries of all exposed sheets at every endpoint contributed by
any other sheet.  Retain the open segments carried by at least two distinct
supporting planes.  Suppress a collinear degree-two point only if its complete
local incident-sheet germ is unchanged through that point.  The resulting
canonical graph is $\GammaM$.  A vertex $a$ \emph{passes} if every neighbour
of $a$ in $\GammaM$ belongs to the perpendicular bisector
\[
 B(o,a)=\{x\in H:\|x-o\|=\|x-a\|\}.
\]
\end{definition}

The definition is intrinsic: exposed sheet subsets, their supporting planes,
and their local germs are properties of $M_J$ as a point set.  In particular,
the graph stores neither a hidden cap label nor an erased internal seam.

A compatible decomposition of $M_J$ consists of a convex core $P'$ and hidden
facets $G_a$ indexed by sites $a$, where the facet set of $P'$ is exactly the
already recovered visible hinge facets together with the pairwise distinct
facets $G_a$, each site $a$ is the reflection of the recovered centre $o$ in
$\aff G_a$, and
\[
 M_J=P'\cup\bigcup_{a\in A}\conv(G_a\cup\{a\})
\]
with pairwise disjoint cell interiors.  The original datum is one such
decomposition.  For an edge $f$ of $G_a$, write
$W(a,f)=\conv(\{a\}\cup f)$ for the corresponding wall.

\begin{lemma}[boundary inventory]
\label{lem:inventory}
For every compatible decomposition,
\[
 \partial M_J
 =\bigcup_{i\in I}F_i\ \cup\
   \bigcup_{a\in A}\ \bigcup_{f\in E(G_a)}W(a,f).
\]
\end{lemma}

\begin{proof}
The boundary of $P'$ is the union of its visible facets and hidden facets.
Each hidden facet is shared with its cap and becomes internal to $M_J$.  The
remaining cap faces are the walls.  Proposition~\ref{prop:contacts} says that
a wall meets another cell only in a lower-dimensional face.
\end{proof}

\begin{lemma}[decomposition computation of the intrinsic graph]
\label{lem:graph-intrinsic}
Starting from the inventory in Lemma~\ref{lem:inventory}, erase a side across
which the two incident boundary patches are coplanar, atomise at every other
sheet-boundary endpoint, and suppress only unchanged collinear degree-two
germs.  The result is $\GammaM$, independently of the compatible
decomposition used for the computation.
\end{lemma}

\begin{proof}
At a relative-interior point of a patch side, exact face-to-face contact gives
no transverse crossing in a patch interior.  If the two incident patch planes
agree, their union is locally planar and the side is absent from the crease
locus.  If the planes differ, the side is an intrinsic noncoplanar crease.
Passing to maximal planar sheets deletes exactly the first kind.  Atomisation
restores every endpoint where another sheet germ begins, ends, or changes,
including a collinear endpoint omitted by a polygon hull.  The final
suppression removes only a presentation subdivision at which the complete
local germ is unchanged.  These conditions depend only on $\partial M_J$.
\end{proof}

\begin{remark}[the swallowed-corner control]
There are legitimate data for which two coplanar merged triangles have a
collinear boundary point that a polygonal hull would drop.  At that point a
noncoplanar incident sheet changes, so Definition~\ref{def:skeleton} retains
it.  The exact control used in the artifact is
\[
 P=\{x\ge0,\ z\ge0,\ x+y\ge2,\ y\le6,\ x+z\le4\},
 \qquad o=(1,2,1),
\]
with the hidden facets $x=0$ and $z=0$.  Their shared edge starts at the foot
of the perpendicular from $o$, making the two apexes and that endpoint
collinear.
\end{remark}

\section{The apex signature and uniqueness}
\label{sec:uniqueness}

\begin{lemma}[only its own walls contain an apex]
\label{lem:own-walls}
For $a\in A$, the boundary patches containing $a$ are exactly the walls of the
cap with apex $a$.
\end{lemma}

\begin{proof}
The reflection $a$ lies strictly outside $P'$.  For $b\ne a$,
Proposition~\ref{prop:contacts} gives
$C_a\cap C_b=G_a\cap G_b\subset P'$, so $a$ belongs to no wall of $C_b$.
Lemma~\ref{lem:inventory} leaves no other boundary patch.
\end{proof}

\begin{lemma}[candidate types]
\label{lem:candidate-types}
Every vertex of $\GammaM$ is either a hidden apex or a vertex of the compatible
core.
\end{lemma}

\begin{proof}
By Lemma~\ref{lem:graph-intrinsic}, a graph vertex is a branch point or a
retained endpoint of the patch-side arrangement.  Exact face-to-face contacts
create no transverse crossing in a patch interior, so such a point is a patch
corner.  Visible-facet corners are core vertices, and wall corners are one
apex and two core vertices.
\end{proof}

\begin{lemma}[apexes pass]
\label{lem:apexes-pass}
Every hidden apex passes in $\GammaM$.
\end{lemma}

\begin{proof}
Let $a$ be the apex over a hidden facet $G_a$.  For each vertex $v$ of $G_a$,
the two walls over the two facet edges incident to $v$ meet along $[a,v]$.
Their planes are distinct: if they were coplanar, the apex and three
consecutive vertices of the convex facet polygon would be coplanar, forcing
$a\in\aff G_a$, contrary to $o\in\operatorname{int}P'$.
Thus the relative interior of $[a,v]$ is an intrinsic crease.  The apex is a
branch point, and $v$ is retained because the two wall-sheet germs containing
$a$ end or change there.  Conversely, Lemma~\ref{lem:own-walls} shows that
every crease incident to $a$ is of this form.  Hence
\[
 N_{\GammaM}(a)=\operatorname{vert}G_a\subset\aff G_a=B(o,a).
\]
\end{proof}

For a facet $f$ of a compatible core, write
$v_f=\vol_3\conv(\{o\}\cup f)$.

\begin{lemma}[volume identity]
\label{lem:volume}
Let $W=\sum_{i\in I}v_{F_i}$.  Every compatible decomposition satisfies
\[
 \vol_3(M_J)=W+2\sum_{a\in A}v_{G_a}.
\]
Consequently all compatible cores have the same volume.
\end{lemma}

\begin{proof}
The radial pyramids from $o$ to the facets tile a convex core, so
$\vol_3(P')=W+\sum_av_{G_a}$.  The cap over $G_a$ is a reflection of its radial
pyramid and has the same volume.  Exact contacts give pairwise disjoint cell
interiors, proving the displayed identity.  Both $M_J$ and $W$ are intrinsic,
so $\sum_av_{G_a}$ and then the core volume are common to all decompositions.
\end{proof}

\begin{lemma}[nested site sets]
\label{lem:nested}
If $(P',A)$ and $(P'',A'')$ are compatible and $A''\subseteq A$, then the two
decompositions coincide.
\end{lemma}

\begin{proof}
The hidden facet associated with $a$ has halfspace
$\{x:\|x-o\|\le\|x-a\|\}$, the Dirichlet-cell description of the core
against its reflected sites~\cite{AshBolker1985}.  Therefore
\[
 P'=\bigcap_{i\in I}H_i^-\cap
     \bigcap_{a\in A}\{x:\|x-o\|\le\|x-a\|\},
\]
and similarly for $P''$.  The inclusion $A''\subseteq A$ gives
$P'\subseteq P''$.  Lemma~\ref{lem:volume} gives equal finite volume, so the
two closed convex bodies agree.  Their active facets and reflected sites then
agree.
\end{proof}

\begin{theorem}[exact reconstruction]
\label{thm:main}
Let $D$ be a three-dimensional proper-zero datum.  Suppose no vertex of its
core passes in $\GammaM$.  Then the passing vertices of $\GammaM$ are exactly
the hidden reflected apexes, the core is
\[
 P=\bigcap_{i\in I}H_i^-\cap
   \bigcap_{a\text{ passing}}\{x:\|x-o\|\le\|x-a\|\},
\]
and every compatible decomposition of $X(D)$ is the same.  Consequently,
$X(D)$ congruent to $X(D')$ implies $D$ congruent to $D'$.
\end{theorem}

\begin{proof}
Let $A$ be the hidden apex set of $D$, and let $(P'',A'')$ be any compatible
decomposition.  The graph $\GammaM$ is intrinsic, hence common to both
decompositions.  Lemma~\ref{lem:apexes-pass} gives that every element of $A$
and $A''$ passes in this common graph.  By Lemma~\ref{lem:candidate-types}, a
passing vertex is an apex of $D$ or a core vertex.  The second possibility is
excluded by hypothesis, so the passing set is exactly $A$ and $A''\subseteq A$.
Lemma~\ref{lem:nested} completes the uniqueness proof.

An ambient isometry carrying one raw union to another preserves the intrinsic
components, centre, exposed sheet germs, and $\GammaM$.  It carries the first
datum to a compatible decomposition of the second raw union, which uniqueness
identifies with the second datum.
\end{proof}

\section{The exceptional centre set}
\label{sec:genericity}

\begin{theorem}[fixed-realization genericity]
\label{thm:genericity}
Fix a realized three-polytope $P$ and a visible-facet mask.  The set of centres
$o\in\operatorname{int}P$ at which some core vertex passes is contained in a
finite union of proper spheres and quadrics.  In particular, it has
three-dimensional Lebesgue measure zero.
\end{theorem}

\begin{proof}
Use fixed absolute coordinates.  Suppose a core vertex $v$ is incident to a
hidden facet with outward unit normal $n$ and equation $n\cdot x=h$.  The
reflected apex is
\[
 a=o+2(h-n\cdot o)n.
\]
The two walls meeting along $[v,a]$ have distinct planes, so this is an edge of
$\GammaM$.  If $v$ passes, then $a\in B(o,v)$, which forces
\[
 4(h-n\cdot o)^2=\|v-o\|^2.
\]
Its quadratic part is $o^{\mathsf T}(4nn^{\mathsf T}-I)o$.  The eigenvalue
along $n$ is $3$ and the two eigenvalues on $n^\perp$ are $-1$.  The polynomial
is nonzero and its zero set is a proper quadric.

If $v$ is incident only to visible hinge facets, the core edges at $v$ are
intersections of distinct facet planes and are intrinsic crease edges.  Any
crease neighbour $w$ required by passing satisfies
\[
 \|w-o\|^2=\|w-v\|^2,
\]
a proper sphere condition in $o$.

There are finitely many vertices and incidences.  Every passing event is
contained in at least one of the displayed proper algebraic hypersurfaces, and
a finite union of such zero sets has measure zero.
\end{proof}

The theorem deliberately does not assert that the bad-centre set itself is
closed.  A subset of a closed algebraic union need not be closed.  Nor does the
argument define a measure on simultaneously varying realizations of $P$.

\section{Ambiguities require five hidden sites per side}
\label{sec:exchange}

This section studies what a failure of uniqueness would have to look like.
Throughout, $D$ is a three-dimensional proper-zero datum with decomposition
$(P,A)$ and centre $o$, and $(P'',A'')$ is a second compatible decomposition
of the same raw union.  By Section~\ref{sec:intrinsic-recovery} the two
decompositions share $M_J$, the centre, the visible hinge facets, and the
intrinsic graph $\GammaM$.

\begin{lemma}[differing sites avoid visible hinges]
\label{lem:diff-site}
Let $a\in A''\setminus A$.  Then $a$ is a vertex of $P$ that passes in
$\GammaM$ and lies on no visible hinge facet.  Symmetrically for
$A\setminus A''$.
\end{lemma}

\begin{proof}
As an apex of the second decomposition, $a$ passes
(Lemma~\ref{lem:apexes-pass}) and, by Lemma~\ref{lem:candidate-types} applied
to the first decomposition, $a$ is an apex of $D$ or a vertex of $P$; the
first option is excluded by $a\notin A$.  Applying
Lemma~\ref{lem:own-walls} in the second decomposition, the only boundary
patches containing $a$ are the walls of its own cap there.  A visible hinge
facet $F_i$ is a boundary patch of every decomposition and is contained in the
core of each; since an apex lies strictly outside the core, $F_i$ is not a
wall of the cap of $a$, and hence $a\notin F_i$.
\end{proof}

\begin{theorem}[hinge-anchored uniqueness]
\label{thm:hinge-anchored}
Suppose every vertex of the core that passes in $\GammaM$ lies on at least one
visible hinge facet.  Then the compatible decomposition is unique.
Theorem~\ref{thm:main} is the special case in which no core vertex passes.
\end{theorem}

\begin{proof}
Suppose $(P'',A'')\neq(P,A)$.  If $A''\subseteq A$,
Lemma~\ref{lem:nested} forces equality, so $A''\setminus A\neq\varnothing$.
Any $a\in A''\setminus A$ is, by Lemma~\ref{lem:diff-site}, a passing core
vertex on no visible hinge facet, contradicting the hypothesis.
\end{proof}

\begin{lemma}[exchange edges]
\label{lem:exchange-edges}
Let $a\in A''\setminus A$.  Then $a$ lies on at least three hidden facets of
$(P,A)$.  For every hidden facet $G_c\ni a$, its site $c$ satisfies
$c\in A\setminus A''$, the segment $[a,c]$ is an edge of $\GammaM$, and the
triangle $(o,a,c)$ is equilateral.  Symmetrically with the two decompositions
exchanged.
\end{lemma}

\begin{proof}
By Lemma~\ref{lem:diff-site}, $a$ is a vertex of the three-polytope $P$, so at
least three facets of $P$ pass through it, and none of them is visible;
distinct facets have distinct planes, and the site is the reflection of $o$ in
the plane, so the corresponding sites are distinct.

Let $G_c$ be such a hidden facet.  Then $a\in\operatorname{vert}G_c$, so by
the proof of Lemma~\ref{lem:apexes-pass} the segment $[c,a]$ is an edge of
$\GammaM$.  The facet plane $\aff G_c$ is the perpendicular bisector of $o$
and $c$, and $a$ lies on it, so $\|a-o\|=\|a-c\|$.  Since $a$ passes and $c$
is one of its graph neighbours, $\|c-o\|=\|c-a\|$.  The triangle is
equilateral.

Finally $c\notin A''$: otherwise $c$ would be an apex of the second
decomposition, and Lemma~\ref{lem:apexes-pass} applied there would give
$N_{\GammaM}(c)=\operatorname{vert}G''_c\subset P''$; but
$a\in N_{\GammaM}(c)$ is an apex of the second decomposition and lies strictly
outside $P''$.
\end{proof}

\begin{lemma}[two common neighbours]
\label{lem:common}
Let $X$ be the graph on $(A\setminus A'')\cup(A''\setminus A)$ whose edges are
the pairs $\{a,c\}$ of Lemma~\ref{lem:exchange-edges}.  Then $X$ is simple and
bipartite between the two difference sets, with minimum degree at least three,
and two distinct sites on one side have at most two common neighbours.
\end{lemma}

\begin{proof}
Bipartiteness, simplicity and the degree bound are
Lemma~\ref{lem:exchange-edges}.  Let $c$ be a common neighbour of the distinct
same-side sites $a$ and $a'$.  The two equilateral relations give
$\|c-o\|=\|c-a\|=\|a-o\|$ and $\|c-o\|=\|c-a'\|=\|a'-o\|$.  So $c$ lies on
the perpendicular bisector plane of $o$ and $a$, on that of $o$ and $a'$, and
on the sphere of radius $\|a-o\|$ about $o$.  The two planes are distinct,
because reflecting $o$ across them returns $a$, respectively $a'$; hence they
meet in a line or not at all, and a line meets a sphere in at most two points.
\end{proof}

\begin{lemma}[incidence count]
\label{lem:count}
Each side of $X$ has at least four sites.  If some side has exactly four,
then both do, $X$ is $3$-regular, that is, $K_{4,4}$ minus a perfect
matching, and all eight sites lie at one common distance $R>0$ from $o$ with
every edge a chord of length $R$.
\end{lemma}

\begin{proof}
Write $p$ and $q$ for the two side sizes.  Every vertex has at least three
neighbours, all on the other side, so $p,q\ge3$.  Count the incidences
$(\{a,a'\},c)$ of a common neighbour $c$ with an unordered pair of distinct
sites on the $p$-side: Lemma~\ref{lem:common} bounds them by
$2\binom p2$, and they number
$\sum_{c}\binom{\deg c}{2}\ge3q$, because each of the $q$ vertices $c$ of the
other side has degree at least three.

If $p=3$, every vertex of the other side is adjacent to all three sites, so
one pair on the $p$-side has $q\ge3>2$ common neighbours, impossible.  Hence,
exchanging the roles of the sides, $p,q\ge4$.  If $p=4$, then
$3q\le2\binom42=12$ forces $q\le4$, hence $q=4$; and a vertex of degree four
on either side would give
$\sum_c\binom{\deg c}{2}\ge\binom42+3\binom32=15>12$, so every degree is
exactly three.  The bipartite complement of a $3$-regular bipartite graph on
$4+4$ vertices is $1$-regular, a perfect matching.  Such a graph is connected
(a component would need at least three vertices on each side), and the
equilateral relation of Lemma~\ref{lem:exchange-edges} propagates the
distance $\|a-o\|$ across every edge, so all eight sites lie at one distance
$R>0$ and every edge has length $R$.
\end{proof}

The remaining case is a genuinely geometric one, and it is the heart of the
bound.

\begin{theorem}[no exchange quadruple]
\label{thm:fourbyfour}
There is no configuration of distinct unit vectors $d_1,\ldots,d_4$ and
distinct unit vectors $e_1,\ldots,e_4$ in $\R^3$ with
\[
 e_l\cdot d_i=\tfrac12
 \qquad\text{for all }l\ne i.
\]
\end{theorem}

\begin{proof}
Write $x_{ij}=d_i\cdot d_j$, let $G$ be the unit-diagonal Gram matrix of
$d_1,\ldots,d_4$, let $G_i$ be its principal $3\times3$ minor deleting row and
column $i$, and let $r_i$ be the deleted off-diagonal column.  Define the
polynomials
\begin{align*}
 g_0&:=\det G,\\
 g_i&:=\mathbf 1^{\mathsf T}\operatorname{adj}(G_i)\mathbf 1-4\det G_i
 \qquad(i=1,\ldots,4),\\
 s_i&:=\mathbf 1^{\mathsf T}\operatorname{adj}(G_i)r_i-\det G_i .
\end{align*}
For example, in the variables $(x_{23},x_{24},x_{34})$,
\begin{multline*}
 g_1=3x_{23}^2+3x_{24}^2+3x_{34}^2-8x_{23}x_{24}x_{34}\\
     +2(x_{23}x_{24}+x_{23}x_{34}+x_{24}x_{34})
     -2(x_{23}+x_{24}+x_{34})-1 .
\end{multline*}

\emph{Step 1: the configuration satisfies the polynomial system.}
Four vectors in $\R^3$ give $g_0=\det G=0$.  Fix $l$ and suppose the triple
$\{d_i:i\ne l\}$ were linearly dependent; applying $e_l\cdot(\,\cdot\,)$ to a
dependence shows the coefficients sum to zero, so three distinct points would
be affinely dependent, hence collinear, and a line meets the unit sphere in at
most two points.  So each triple is a basis and $\det G_l\neq0$, and $e_l$ is
the unique solution of the three linear equations $e_l\cdot d_i=\tfrac12$
($i\ne l$).  Solving that system by Cramer's rule and imposing
$\|e_l\|^2=1$ yields
$\mathbf 1^{\mathsf T}\operatorname{adj}(G_l)\mathbf 1=4\det G_l$, that is,
$g_l=0$ on the Gram coordinates of the configuration.  Moreover
$c_l:=e_l\cdot d_l-\tfrac12$ satisfies $2c_l\det G_l=s_l$.  If some $c_l=0$,
then $e_l\cdot d_i=\tfrac12$ for all four $i$, so all four $d_i$ lie on the
plane $\{x:e_l\cdot x=\tfrac12\}$, which misses the origin; any three of them
are then linearly independent, and every $e_m$ solves the same three linear
equations as $e_l$, forcing $e_m=e_l$ for all $m$ and contradicting
distinctness.  Hence every $s_l\neq0$.

\emph{Step 2: four certificate identities force an isosceles pattern.}
There exist explicit polynomials
$q_0^{(\ell)},\ldots,q_4^{(\ell)}\in\mathbb Q[x_{12},\ldots,x_{34}]$ of total
degree at most eight with
\[
 \ell\cdot s_1s_2
 \;=\;
 q_0^{(\ell)}g_0+q_1^{(\ell)}g_1+q_2^{(\ell)}g_2+q_3^{(\ell)}g_3
 +q_4^{(\ell)}g_4
\]
for each of the four linear forms
\[
 \ell\in\{\,x_{12}-x_{34},\;x_{13}-x_{24},\;x_{14}-x_{23},\;
          x_{23}+x_{24}+x_{34}+1\,\}.
\]
The cofactors are archived, with the verification script, alongside the
reference implementation; each identity is re-verified by expanding the
difference of the two sides to the zero polynomial in exact rational
arithmetic, independently of the software that found it.  At the configuration all
$g_i$ vanish and $s_1s_2\neq0$, so all four linear forms vanish:
\[
 x_{12}=x_{34}=:a,\qquad
 x_{13}=x_{24}=:b,\qquad
 x_{14}=x_{23}=:c,\qquad
 a+b+c=-1 .
\]

\emph{Step 3: an elementary emptiness argument.}
Substituting the pattern into any of the four face polynomials leaves one
cubic; explicitly,
\[
 g_i\big|_{\text{pattern}}=4F(a,b),
 \qquad
 F(a,b)=2ab^2+2a^2b+a^2+3ab+b^2+a+b+1,
\]
an identity in $(a,b)$ that is checked by direct substitution.  Hence
$F(a,b)=0$.  By Cauchy--Schwarz $|x_{ij}|\le1$, so $(a,b)\in[-1,1]^2$.  On
the open square $F$ is strictly positive.  Indeed the difference
$\partial_bF-\partial_aF$ factors as $(a-b)(2a+2b+1)$; on the branch $b=a$
the equation $\partial_aF=0$ becomes $6a^2+5a+1=(3a+1)(2a+1)=0$, and on the
branch $2a+2b+1=0$ it becomes $-a(2a+1)=0$.  This gives exactly the four
interior critical points
$(-\tfrac13,-\tfrac13),(-\tfrac12,-\tfrac12),(0,-\tfrac12),(-\tfrac12,0)$
with values $\tfrac{20}{27},\tfrac34,\tfrac34,\tfrac34$, and the edge
restrictions of $F$ are
\[
\begin{aligned}
 F(a,1)&=3(a+1)^2, &\qquad F(a,-1)&=1-a^2,\\
 F(1,b)&=3(b+1)^2, &\qquad F(-1,b)&=1-b^2,
\end{aligned}
\]
whose only zeros on the closed square are the three corners $(-1,1)$,
$(1,-1)$ and $(-1,-1)$ of the $(a,b)$-square.  An interior zero of $F$ would
force an interior critical point with a nonpositive value, and there is none.
At the three corners, respectively $x_{24}=1$, $x_{34}=1$, and
$x_{14}=-1-a-b=1$, so two of the $d_i$ coincide by the equality case of
Cauchy--Schwarz, contradicting distinctness.  No configuration exists.
\end{proof}

\begin{remark}
The certificate identities show that the four linear forms, and with them
$F(x_{34},x_{24})$, lie in the saturation of
$\langle g_0,\ldots,g_4\rangle$ by $s_1s_2$: over the complex numbers, every
solution of the distance-$\tfrac12$ system with $s_1s_2\neq0$ is already
isosceles-patterned.  The certificates were found by a Gr\"obner-basis
saturation in Singular~4.3.2, with the saturation recomputed by iterated
ideal quotients under in-run containment and stability checks, and the
computed reduced basis of the fully saturated ideal is
$\{x_{23}+x_{24}+x_{34}+1,\ x_{14}+x_{24}+x_{34}+1,\ x_{13}-x_{24},\
x_{12}-x_{34},\ F(x_{34},x_{24})\}$; the proof above uses only the certified
fragment, re-verified by exact expansion in SymPy~1.14, so no Gr\"obner
implementation is trusted.  As a standing positive control, the same pipeline
at the distance $\tfrac13$ keeps the regular tetrahedron on the saturated
variety, as it must.
\end{remark}

\begin{theorem}[five sites per side]
\label{thm:five}
Two distinct compatible decompositions of one proper-zero raw union differ in
at least five reflected sites on each side:
$|A\setminus A''|\ge5$ and $|A''\setminus A|\ge5$.
\end{theorem}

\begin{proof}
As in Theorem~\ref{thm:hinge-anchored}, distinctness makes both difference
sets nonempty.  By Lemma~\ref{lem:count} each side of the exchange graph $X$
has at least four sites, and if some side had exactly four, $X$ would be
$K_{4,4}$ minus a perfect matching with all eight sites at one distance
$R>0$ from $o$ and all twelve edges of length $R$.  Translating $o$ to the
origin and scaling by $R^{-1}$ makes the sites unit vectors with
$\|a-c\|=1$, i.e.\ $a\cdot c=\tfrac12$, along every edge; labelling the
matching pairs with equal indices produces exactly the configuration of
Theorem~\ref{thm:fourbyfour}, which does not exist.  Hence both sides have at
least five sites.
\end{proof}

\begin{corollary}[at most four hidden facets]
\label{cor:four}
A proper-zero datum with at most four hidden facets has a unique compatible
decomposition, with no hypothesis on its centre.
\end{corollary}

\begin{proof}
A second decomposition would give $|A\setminus A''|\ge5$, while
$|A|\le4$.
\end{proof}

\section{Exact reconstruction algorithm and evidence}
\label{sec:algorithm}

Assume an exact rational polygonal presentation of the exposed boundary.
Coplanar subdivisions are presentation data and may be inserted or removed.
The conceptual algorithm is:

\begin{enumerate}[label=\arabic*.]
\item Compute the closures of the connected components of the local-flat
locus.  With at least three components, select the unique centre of their
codimension-one adjacency star; with exactly two, select the component of
larger intrinsic volume.  This is $M_J$.
\item Intersect every other component with $M_J$ to recover the nonflat hinges
and their cap geometry.
\item Recover $o$ from either nonflat cap by the two-candidate construction in
Proposition~\ref{prop:intrinsic-recovery}.
\item Construct $\GammaM$ from exposed sheet germs.
\item Keep exactly the graph vertices that pass.
\item Intersect the visible hinge halfspaces with the bisector halfspaces of
the passing vertices, the Dirichlet-cell construction
of~\cite{AshBolker1985}.
\end{enumerate}

Under the hypothesis of Theorem~\ref{thm:main}, correctness is immediate from
Lemmas~\ref{lem:candidate-types} and~\ref{lem:apexes-pass}.  A full certifier
may rebuild the entire raw union, including the observed nonflat caps, and
compare it exactly with the input.  The reference code implements the flat
reconstruction kernel in steps 4--6; it is not an end-to-end implementation of
steps 1--3 or of the full raw-union certifier.

The certifying variant compares two intrinsic conditions, and each is provably
necessary.  There are constructed exact data, at exceptional centres, on
which the reconstruction returns a strict superset of the true site set while
the rebuilt merged flat component still equals the observed one exactly:
cutting a passing vertex off with its bisector plane and re-attaching the cap
over the new facet restores the removed piece.  Comparing only the merged flat
component therefore certifies a wrong answer on those data (six exact
witnesses).  What the wrong reconstruction destroys is the recovered visible
hinge facets, so the certifier also requires every recovered visible hinge
facet to survive as a facet of the reconstructed core.  On a second witness
family (three exact witnesses) the situation is reversed: the hinge facets
survive while the flat components differ.  Hence neither condition alone
suffices, and the certifier checks both.

Let $N$ be the number of input boundary corners and $k$ the number of returned
core facets.  The reference germ extractor atomises each segment against all
exact endpoints and builds carrier labels during atomisation.  Its conservative
bound is $O(N^2\log N)$ exact field operations and $O(N^2)$ crease atoms.  The
reference three-dimensional halfspace kernel enumerates the $\binom k3$ plane
triples and validates each candidate corner against every halfspace, so it
costs $O(k^4)$.  Thus the proved reference bound is
\[
 O(N^2\log N+k^4).
\]
No sharper literature bound is imported into this statement.

The executable audit has two layers.  E-A60 runs the intrinsic germ skeleton
on every frozen six-vertex row.  E-A58 runs the reconstruction/flat
certification kernel.  The frozen measurements are:

\begin{center}
\begin{tabular}{lr}
\hline
intrinsic-skeleton corpus rows & 16,744\\
rows reconstructed by the intrinsic signature & 16,744\\
rows agreeing with the historical graph on the corpus & 16,744\\
subdivision mutations invariant (256-row corpus prefix) & 256/256\\
maximal-sheet cross-checks (same prefix) & 256/256\\
flat-kernel reconstructions & 16,744/16,744\\
sampled two-condition certifications & 340/340\\
wrong answers certified & 0\\
degenerate swallowed-corner witness & pass\\
exceptional-centre witnesses & 6\\
flat-only accepts / two-condition rejects & 6/6\\
reversed witnesses (hinges survive, flat differs) & 3/3\\
corpus disagreements between the two certifier variants & 0/340\\
certificate identities of Theorem~\ref{thm:fourbyfour} verified & 4/4\\
\hline
\end{tabular}
\end{center}

The canonical payloads are, for the intrinsic-skeleton audit E-A60,
\begin{center}
\small\texttt{D18C07176BAF7B729A550E9B9B8B9D9E04BB58A0C52D295243108417C9BE785A},
\end{center}
for the two-condition reconstruction/certification kernel E-A58,
\begin{center}
\small\texttt{3C58C0B88081380C7D22993A6FEE3B4EC9FE1A8C3D0CDC107BE9D40C33BA4ACC},
\end{center}
for the exceptional-centre adversarial control E-A61,
\begin{center}
\small\texttt{76956805B0BA304142CB08D833FA1179876A9F2154A76CE6C5A232BEF9E3F7F1},
\end{center}
for the reversed witness family E-A63,
\begin{center}
\small\texttt{50B26633575C81C67AD3A9857DC5BE1227B512E21CC98168EECEC37B360DB1A8},
\end{center}
and for the certificate verification of Theorem~\ref{thm:fourbyfour}, whose
payload pins the SHA-256 of all four cofactor files, E-A69,
\begin{center}
\small\texttt{5DC4C5800572C64F07C7A815CE7BEFFB930FEF03E04CC71679FB9501606F1D92}.
\end{center}
Finite replay is evidence about the implementation and its frozen domain.  The
proof of Theorem~\ref{thm:main} is the intrinsic argument in the preceding
sections, and the proof of Theorem~\ref{thm:fourbyfour} rests on the verified
identities and the elementary analysis of Section~\ref{sec:exchange}, not on
the replay.

\section{Relation to supplied-boundary inverse problems}
\label{sec:related}

Ash and Bolker study recognition of supplied Dirichlet tessellations.  Their
Lemma~1 places a common cell boundary on the perpendicular bisector of its two
sources, and their Theorem~2 constructs a Dirichlet cell by reflecting an
interior point in the known facet hyperplanes~\cite{AshBolker1985}.  The latter
is exactly the forward reflection construction used in Section~\ref{sec:model}
and is cited as such.

Subsequent recognition algorithms also begin with substantially visible
boundary data.  Aurenhammer takes a supplied polytopal cell
complex~\cite{Aurenhammer1987}; Hartvigsen takes a tessellation of Euclidean
space into polyhedra and uses linear programming~\cite{Hartvigsen1992}; and
Schoenberg, Ferguson, and Li recover generators from segments demarcating cell
boundaries~\cite{SchoenbergFergusonLi2003}.  In the plane, Alonso Ferrero states
the generator-recognition problem with the Voronoi edges
provided~\cite{AlonsoFerrero2011}.  Biedl, Held, and Huber reconstruct from a
supplied embedded straight-line graph with rays~\cite{BiedlHeldHuber2013}.
Power-diagram detection and weighted bisector-graph reconstruction similarly
start from a supplied complex, partition, or geometric graph
\cite{BorgwardtFrongillo2017,EderEtAl2023}.  The generalized inverse Voronoi
problem studied by Aloupis et al. permits added edges while retaining the
original input edges~\cite{AloupisEtAl2013}.

Not every inverse tessellation problem receives boundary edges.  Bourne,
Pearce, and Roper prove that labelled Laguerre cells are uniquely determined
by their volumes and centroids and recover them by convex
optimization~\cite{BournePearceRoper2025}.  Bourne, Mulholland, Sahu, and
Tant instead fit an oriented planar Voronoi diagram from ultrasonic boundary
travel times in a simplified nondestructive-testing model; the geometry is
unknown, their first objective can have nonunique minimizers, and the reported
reconstruction is noise-sensitive~\cite{BourneEtAl2021NDT}.  These are
partial- and indirect-data comparators, but their observables are respectively
global labelled cell moments and travel times, not a reflected-cap raw union.

Three further classical theories are sometimes suggested as antecedents and
are not.  Alexandrov's theory determines a convex polyhedron from its
intrinsic surface metric, which is blind to the fold structure of a flat
doubly covered region --- exactly the information at issue
here~\cite{Alexandrov2005}.  The Blind--Mani theorem, with Kalai's proof,
reconstructs a simple polytope from its graph, with combinatorial input and
output~\cite{BlindMani1987,Kalai1988}.  Geometric tomography reconstructs
bodies from projections or sections, and the reflected-cap raw union is
neither~\cite{Gardner2006}.

The precise distinction in the present model is that the zero-angle source
layer and its facet boundaries have already been erased by the raw union.  The
intrinsic germ skeleton is designed to use only the boundary information that
survives that erasure.  This comparison is a statement about input hypotheses,
not a claim that the present result is the first possible treatment of every
related erased-boundary problem.

\section{Limitations and open problems}
\label{sec:limitations}

The following questions are not resolved here.

\begin{itemize}
\item Whether an exact ambiguity exists at all remains open.
Theorems~\ref{thm:hinge-anchored} and~\ref{thm:five} constrain it: every
ambiguous datum needs a passing core vertex avoiding all visible hinge facets
and at least five differing sites on each side, hence at least five hidden
facets.  Data with five or more hidden facets are not decided here.
\item The proof is three-dimensional.  Higher-dimensional analogues require a
new analysis of the boundary-germ strata and candidate types.
\item Theorem~\ref{thm:fourbyfour} is computer-assisted: the four certificate
identities were found by machine and are re-verified by exact expansion.  The
verification chain ships with the reference implementation; the remaining
steps are hand-checkable.
\item The reference implementation assumes an exact polygonal presentation.
It does not extract such a presentation from a point cloud, floating-point
mesh, or arbitrary implicit-set oracle.
\item E-A58 implements the flat reconstruction kernel and the two-condition
certifier; it is not the full six-step pipeline or a raw-union certifier.
Completeness of the two conditions is not claimed; each is necessary, and
their conjunction rejects every constructed wrong reconstruction.
\item The finite corpus and deterministic mutation sample do not establish
novelty, priority, or unrestricted-domain behaviour.
\end{itemize}

\begingroup
\small
\bibliographystyle{plain}
\bibliography{references}
\endgroup

\end{document}
